% GENERATED FILE. Edit canonical lessons, then run npm run generate:books.
% canonical-source-hash: fnv1a64:fd9889f42d70591c
% canonical-lessons: IT-C344-prova, IT-C344-evidente, IT-C344-ovvio, IT-C344-indiscutibile, IT-C344-palese

\chapter{Le prove}
\label{ch:it-le-prove}
\hlchaptermodality{344}

\begin{chapteropening}
\textbf{By the end of this chapter:} \emph{I can talk about proof: say something is evident, obvious or beyond question, or that it is plain to see.}

Complete the last lesson of chapter 344: I can talk about proof: say something is evident, obvious or beyond question, or that it is plain to see.
\end{chapteropening}

\section[la prova]{la prova — the proof, the evidence; the test, the trial}

\label{lesson:IT-C344-prova}



\begin{quote}
Before the new one: say the Italian for absurd, then the Italian for questionable.
\end{quote}

\subsection*{You'll want to know: la prova}
\textbf{la prova} — \textquotedblleft{}the proof, the evidence; the test, the trial\textquotedblright{}.

\begin{quote}
Non ci sono prove che sia stato lui.
\end{quote}

\begin{quote}
\emph{There is no proof that it was him.}
\end{quote}

\emph{Mettere alla prova} — to put to the test. In a shop, \emph{la prova} is also trying something on.

\begin{cousinweb}
\emph{prova} is built on \emph{provare}, inherited from Latin \emph{probāre} \textquotedblleft{}to test, to approve, to prove\textquotedblright{} (Lewis \& Short), from \emph{probus} \textquotedblleft{}good, upright\textquotedblright{}.
\end{cousinweb}

\begin{grammarlens}[title={What you've built}]
One more word for telling, explaining and describing at B1.
\end{grammarlens}

\subsection*{Your turn}
\begin{itemize}
\raggedright
  \item \emph{Say it:} \emph{la prova}
  \item \emph{Say it:} \emph{la prova}, once more
  \item \emph{Say it:} a sentence of your own with \emph{la prova}
  \item \emph{Recall:} say the Italian for absurd, then the Italian for questionable, then say \emph{la prova} again
\end{itemize}

\begin{tcolorbox}[breakable,colback=teal!4,colframe=teal!35!black,title={Before you move on}]
What does la prova mean? (\textquotedblleft{}The proof, the evidence; the test, the trial\textquotedblright{}.) Say it once more.
\end{tcolorbox}
\section[evidente]{evidente — evident, clear}

\label{lesson:IT-C344-evidente}



\begin{quote}
Before the new one: say the Italian for questionable, then the Italian for the proof.
\end{quote}

\subsection*{You'll want to know: evidente}
\textbf{evidente} — \textquotedblleft{}evident, clear\textquotedblright{}.

\begin{quote}
È evidente che non ha studiato.
\end{quote}

\begin{quote}
\emph{It is clear that he has not studied.}
\end{quote}

\emph{È evidente che} takes the indicative, because the fact is certain.

\begin{cousinweb}
\emph{evidente} is a learned borrowing from Latin \emph{ēvidēns} \textquotedblleft{}visible, apparent, clear\textquotedblright{} (EtymDB), built from \emph{ē} \textquotedblleft{}out\textquotedblright{} and \emph{vidēre} \textquotedblleft{}to see\textquotedblright{}: something that stands out to the eye.
\end{cousinweb}

\begin{grammarlens}[title={What you've built}]
One more word for telling, explaining and describing at B1.
\end{grammarlens}

\subsection*{Your turn}
\begin{itemize}
\raggedright
  \item \emph{Say it:} \emph{evidente}
  \item \emph{Say it:} \emph{evidente}, once more
  \item \emph{Say it:} a sentence of your own with \emph{evidente}
  \item \emph{Recall:} say the Italian for questionable, then the Italian for the proof, then say \emph{evidente} again
\end{itemize}

\begin{tcolorbox}[breakable,colback=teal!4,colframe=teal!35!black,title={Before you move on}]
What does evidente mean? (\textquotedblleft{}Evident, clear\textquotedblright{}.) Say it once more.
\end{tcolorbox}
\section[ovvio]{ovvio — obvious}

\label{lesson:IT-C344-ovvio}



\begin{quote}
Before the new one: say the Italian for the proof, then the Italian for evident.
\end{quote}

\subsection*{You'll want to know: ovvio}
\textbf{ovvio} — \textquotedblleft{}obvious\textquotedblright{}.

\begin{quote}
Mi sembra ovvio: se piove, restiamo a casa.
\end{quote}

\begin{quote}
\emph{It seems obvious to me: if it rains, we stay at home.}
\end{quote}

\emph{Ovviamente} — obviously. Plural \emph{ovvi}, feminine \emph{ovvia}.

\begin{cousinweb}
\emph{ovvio} is a learned borrowing from Latin \emph{obvius} \textquotedblleft{}in the way, meeting, easy to reach\textquotedblright{} (Lewis \& Short), from \emph{ob viam} \textquotedblleft{}towards the road\textquotedblright{}: what you meet on the path.
\end{cousinweb}

\begin{grammarlens}[title={What you've built}]
One more word for telling, explaining and describing at B1.
\end{grammarlens}

\subsection*{Your turn}
\begin{itemize}
\raggedright
  \item \emph{Say it:} \emph{ovvio}
  \item \emph{Say it:} \emph{ovvio}, once more
  \item \emph{Say it:} a sentence of your own with \emph{ovvio}
  \item \emph{Recall:} say the Italian for the proof, then the Italian for evident, then say \emph{ovvio} again
\end{itemize}

\begin{tcolorbox}[breakable,colback=teal!4,colframe=teal!35!black,title={Before you move on}]
What does ovvio mean? (\textquotedblleft{}Obvious\textquotedblright{}.) Say it once more.
\end{tcolorbox}
\section[indiscutibile]{indiscutibile — unquestionable, beyond question}

\label{lesson:IT-C344-indiscutibile}



\begin{quote}
Before the new one: say the Italian for evident, then the Italian for obvious.
\end{quote}

\subsection*{You'll want to know: indiscutibile}
\textbf{indiscutibile} — \textquotedblleft{}unquestionable, beyond question\textquotedblright{}.

\begin{quote}
La sua bravura è indiscutibile.
\end{quote}

\begin{quote}
\emph{Her skill is beyond question.}
\end{quote}

Built on \emph{discutibile} with \emph{in-} \textquotedblleft{}not\textquotedblright{}.

\begin{cousinweb}
\emph{indiscutibile} is built from \emph{in-} \textquotedblleft{}not\textquotedblright{} and \emph{discutibile}, on \emph{discutere} \textquotedblleft{}to discuss\textquotedblright{}, from Latin \emph{discutere} \textquotedblleft{}to strike apart\textquotedblright{} (Lewis \& Short), built on \emph{quatere} \textquotedblleft{}to shake\textquotedblright{}.
\end{cousinweb}

\begin{grammarlens}[title={What you've built}]
One more word for telling, explaining and describing at B1.
\end{grammarlens}

\subsection*{Your turn}
\begin{itemize}
\raggedright
  \item \emph{Say it:} \emph{indiscutibile}
  \item \emph{Say it:} \emph{indiscutibile}, once more
  \item \emph{Say it:} a sentence of your own with \emph{indiscutibile}
  \item \emph{Recall:} say the Italian for evident, then the Italian for obvious, then say \emph{indiscutibile} again
\end{itemize}

\begin{tcolorbox}[breakable,colback=teal!4,colframe=teal!35!black,title={Before you move on}]
What does indiscutibile mean? (\textquotedblleft{}Unquestionable, beyond question\textquotedblright{}.) Say it once more.
\end{tcolorbox}
\section[palese]{palese — plain, open (for all to see)}

\label{lesson:IT-C344-palese}



\begin{quote}
Before the new one: say the Italian for obvious, then the Italian for unquestionable.
\end{quote}

\subsection*{You'll want to know: palese}
\textbf{palese} — \textquotedblleft{}plain, open (for all to see)\textquotedblright{}.

\begin{quote}
La sua gelosia era palese.
\end{quote}

\begin{quote}
\emph{His jealousy was plain to see.}
\end{quote}

Formal; in speech \emph{evidente} or \emph{chiaro} are more frequent.

\begin{cousinweb}
\emph{palese} is usually traced to Latin \emph{pālam} \textquotedblleft{}openly, in public\textquotedblright{} (Lewis \& Short), through a Late Latin adjective of the type \emph{palēnsis}; the exact path is not settled.
\end{cousinweb}

\begin{grammarlens}[title={What you've built}]
One more word for telling, explaining and describing at B1.
\end{grammarlens}

\subsection*{Your turn}
\begin{itemize}
\raggedright
  \item \emph{Say it:} \emph{palese}
  \item \emph{Say it:} \emph{palese}, once more
  \item \emph{Say it:} a sentence of your own with \emph{palese}
  \item \emph{Recall:} say the Italian for obvious, then the Italian for unquestionable, then say \emph{palese} again
\end{itemize}

\begin{tcolorbox}[breakable,colback=teal!4,colframe=teal!35!black,title={Before you move on}]
What does palese mean? (\textquotedblleft{}Plain, open (for all to see)\textquotedblright{}.) Say it once more.
\end{tcolorbox}
